\section{Summary and Future Work}

This paper develops a Euclidean, \emph{Principia}-style treatment of the Kepler problem using finite-step constructions, tangent/triangle geometry, and Newton's ultimate-ratio viewpoint. For the direct problem, the proof chain across Sections~3--7 is now explicit and modular. Section~3 establishes the dynamical bridge: from constant areal speed, the acceleration direction is central, and the two local lemmas convert a local geometric ratio into the inverse-square force form. Section~4 sets the geometric transport layer by introducing the auxiliary circle and affine map, so tangent and parallel structures can be moved cleanly between circle and ellipse. Section~5 executes the core local-drop computation in this transported geometry: tangent transfer, matching normal drops, and swept-area transport reduce the local deflection to a conic-invariant ratio, and Section~5.7 is stated as Proof~1 for the ellipse. The same architecture is used throughout: conic-specific geometry computes the local ratio, while the Section~3 conversion step remains universal. Section~6 then gives Proof~2 and Proof~3 as independent auxiliary-circle derivations, confirming the same force law from different local identity chains and clarifying the hodograph-facing interpretation. Section~7 applies the same architecture to the non-elliptic conics: hyperbola and parabola are treated as geometry replacements in the local-ratio step, with the same final conversion to the inverse-square law.

For the reverse problem (Section~8), we keep the hodograph-circle theorem in view but work in a discrete Newtonian style. The key construction is the translated hodograph scaled by $\Delta t$ and carried along the orbit: previous, current, and next velocities are placed together in the same small-circle picture, and the corresponding tangent/offset geometry is used to reconstruct the conic form in configuration space. This gives a diagrammatic bridge from velocity-space circularity to orbit-shape recovery without relying on a continuous differential-equation derivation at each step. The main limitation remains uniformity of local geometry across conic types: the elliptic affine-circle method is especially clean, while parabola and hyperbola require additional case-specific constructions. The next step is therefore to package Sections~5--7 into a single conic-uniform local-drop framework, so that one canonical geometric template handles all three energy regimes while preserving the same Section~3 conversion principle and the same translated, $\Delta t$-scaled hodograph interpretation in the reverse direction.
